% Esquisse d'un programme (cote 119, copy 2): archive cote 119, site page 90 (PDF page 91), typescript page 12.
% typescript. Draft compared: esquisse/pages/esq_119_014.tex (copy 1, site 14).
% AI-transcribed by Claude Opus 5.5 (claude-opus-5-5, adaptive thinking, effort high), Message Batches API, batch esq119_full (msgbatch_015xqDHkTumDysRF1qngZhRJ), 2026-11-13; not yet checked by eye.
\scanpage{119/90}
\mspage{12}

projective complexe est définie sur le corps de base absolue $\mathbb{Q}$ ,
ainsi que les points de ramification admis, les courbes algébriques
obtenues sont définies non seulement sur $\mathbb{C}$, mais sur la clôture
algébrique $\overline{\mathbb{Q}}$ de $\mathbb{Q}$ dans $\mathbb{C}$. Quant à la carte de départ, elle se
retrouve sur la courbe algébrique, comme image inverse du segment
réel $[0,1]$ (où $0$ est considéré comme un sommet, et $1$ comme
milieu d'une ``arête pliée'' ayant 1 comme centre), lequel constitue
dans la sphère de Riemann la ``2-carte orientée universelle'' (*).
Les points de la courbe algébrique $X$ au-dessus de 0, de 1 et de \unsure{\hand{$\infty$}|$\infty$}
ne sont autres que les sommets, et les ``centres'' des arêtes et des
faces respectivement de la carte $(X,K)$, et les ordres des sommets
et des faces ne sont autres que les multiplicités des zéros et
des pôles de la fonction rationnelle (définie sur $\overline{\mathbb{Q}}$) sur $X$,
exprimant sa projection structurale vers $\mathbb{P}^1_{\mathbb{C}}$.

Cette découverte, qui techniquement se réduit à si peu
de choses, a fait sur moi une impression très forte, et elle repré-%
sente un tournant décisif dans le cours de mes réflexions, un dé-%
placement notamment de mon centre d'intérêt en mathématique, qui
soudain s'est \inserted{re}trouvé fortement localisé. Je ne crois pas qu'un
fait mathématique m'ait jamais autant frappé que celui-là, et ait
eu un impact psychologique comparable ($^2$). Cela tient sûrement à la
nature tellement familière, non technique, des objets considérés,
dont tout dessin d'enfant griffonné sur un bout de papier (pour peu
que le graphisme soit d'un seul tenant) donne un exemple parfaitement
explicite. A un tel dessin se trouvent associés des invariants arith-%
métiques subtils, qui seront chamboulés complètement dès qu'on y
rajoute un trait de plus. S'agissant ici de cartes sphériques, don-%

\noindent\rule{4cm}{0.4pt}

(*) Il y a une description analogue des cartes finies non orientées,
éventuellement avec bord, en termes de courbes algébriques \uline{réelles},
plus précisément de revêtement de $\mathbb{P}^1_{\mathbb{R}}$ ramifié seulement on 0, 1, $\infty$,
la surface à bord associée à un tel revêtement étant $X(\mathbb{C})/\tau$, où $\tau$
est la conjugaison complexe. La carte non orientée ``universelle'' est
ici le disque, ou hémisphère supérieur de la sphère de Riemann, muni
comme précédemment du 1-complexe plongé $K = [0,1]$.

\margin{90}
